% !TEX root = ../../main.tex

\section{Preliminaries}\label{sec:p2-prelim}

Throughout this paper, $(M^{n+1}, g)$ denotes a smooth, closed (compact without boundary) Riemannian manifold with $n \geq 2$. We use $\mathcal{Z}_{n}(M; \mathbb{Z}_2)$ to denote the space of mod~$2$ flat cycles in $M$, equipped with the flat norm $\mathcal{F}$.

\begin{definition}[Caccioppoli set~{\cite[\S14, p.\,71]{Sim84}; \cite[Chapter~12, p.\,117]{Mag12}; \cite[\S1]{CLS22}}]\label{def:p2-caccioppoli}
A measurable set $E \subset M$ is a \emph{Caccioppoli set} (set of finite perimeter) if
\[
\mathrm{Per}(E) := \sup\left\{ \int_M \chi_E \operatorname{div} \omega \,:\, \omega \in \mathfrak{X}(M),\; \|\omega\|_{C^0} \leq 1 \right\} < \infty,
\]
where $\chi_E$ denotes the indicator function of $E$. The \emph{perimeter of $E$ in an open set $U$}, denoted $\mathrm{Per}(E, U)$, is the total variation of $D\chi_E$ in $U$.
\end{definition}

By De~Giorgi's structure theorem~\cite[Theorem~14.3]{Sim84}, the distributional derivative $D\chi_E$ of a Caccioppoli set $E$ is given by $D\chi_E = \nu_E \, \mathcal{H}^n \mres \partial^* E$, where $\partial^* E$ is the \emph{reduced boundary} of $E$ (an $n$-rectifiable set), and $\nu_E$ is the outward unit normal defined $\mathcal{H}^n$-a.e.\ on $\partial^* E$. This allows us to identify $D\chi_E$ with an element of $\mathcal{Z}_n(M; \mathbb{Z}_2)$, which we denote by $\partial E := \nu_E \, \mathcal{H}^n \mres \partial^* E$. With this identification,
\[
\mathbf{M}(\partial E) = \mathrm{Per}(E) \qquad \text{and} \qquad \mathcal{F}(\partial E, \partial F) = \|\chi_E - \chi_F\|_{L^1} = \mathrm{Vol}(E \triangle F).
\]

A fundamental tool for Caccioppoli sets is the relative isoperimetric inequality, which controls volume by perimeter inside geodesic balls:

\begin{proposition}[Relative isoperimetric inequality~{\cite[Proposition~12.37]{Mag12}}]\label{prop:relative-isoperimetric}
Let $(M^{n+1}, g)$ be a closed Riemannian manifold. There exist constants $r_1 > 0$ and $C_I > 0$ depending on $(M, g, n)$ such that for every Caccioppoli set $\Omega \subset M$, every $p \in M$, and every $\rho < r_1$,
\begin{equation}\label{eq:rel-isop}
    \min\bigl(\mathcal{L}^{n+1}(\Omega \cap B_\rho(p)),\; \mathcal{L}^{n+1}(B_\rho(p) \setminus \Omega)\bigr)^{n/(n+1)} \leq C_I\, \mathrm{Per}(\Omega, B_\rho(p)).
\end{equation}
\end{proposition}

We will occasionally need the language of integral currents, which we briefly recall; see~\cite[\S26--27]{Sim83} or~\cite[Chapter~4]{Fed69} for a full treatment.

\begin{definition}[Integral currents and cycles~{\cite[\S26, p.\,131]{Sim83}}]\label{def:p2-integral-current}
A $k$-dimensional \emph{integral current} $T$ in $\mathbb{R}^{n+1}$ (or in a Riemannian manifold $(M,g)$) is a linear functional on smooth $k$-forms that can be represented by integration over a $k$-rectifiable set with integer-valued multiplicity.
The \emph{mass} of $T$ is $\mathbf{M}(T) := \sup\{T(\omega) : \|\omega\|_{C^0} \leq 1\}$, and the \emph{boundary} $\partial T$ is the $(k-1)$-current defined by $\partial T(\omega) = T(d\omega)$.
An integral current $T$ with $\partial T = 0$ is called an \emph{integral cycle}.
In the codimension-one case ($k = n$ in $\mathbb{R}^{n+1}$), the boundary of a Caccioppoli set $\partial\llrr{E}$ is an $n$-dimensional integral current with $\mathbf{M}(\partial\llrr{E}) = \mathrm{Per}(E)$.
\end{definition}

The min-max construction produces limits not in the space of Caccioppoli sets but in the space of varifolds, which we now recall.

\begin{definition}[Varifold~{\cite[\S38, p.\,228]{Sim84}}]\label{def:p2-varifold}
An \emph{$n$-varifold} on $M$ is a Radon measure on the Grassmannian bundle $\mathrm{Gr}_n(M)$. We denote the space of $n$-varifolds by $\mathcal{V}_n(M)$. Given a cycle $\Gamma = \partial \Omega \in \mathcal{Z}_n(M; \mathbb{Z}_2)$, we write $|\Gamma|$ for the associated integral varifold:
\[
|\Gamma| = \mathcal{H}^n \mres \partial^* \Omega \otimes \delta_{T_x(\partial^* \Omega)}, \qquad \mu_\Gamma := \mathcal{H}^n \mres \partial^* \Omega.
\]
\end{definition}

We will switch freely between the equivalent notations $\mathbf{M}(\Phi(x))$ and $\mathrm{Per}(\Omega(x))$.

For an $\mathcal{H}^n$-rectifiable set $\Sigma$, the approximate tangent plane $T_p\Sigma$ exists and is unique at $\mathcal{H}^n$-a.e.\ point $p \in \Sigma$ (cf.~\cite[Theorem~11.8]{Sim84}).

\begin{definition}[Rectifiable varifold~{\cite[\S15, p.\,77]{Sim84}}]\label{def:rectifiable-varifold}
An $n$-varifold $V$ on $(M^{n+1}, g)$ is \emph{rectifiable} if there exist an $\mathcal{H}^n$-rectifiable set $\Sigma \subset M$ and a locally $\mathcal{H}^n$-integrable function $\theta \colon \Sigma \to (0, \infty)$ such that $V = \underline{v}(\Sigma, \theta)$, i.e.,
\[
\int f \, dV = \int_\Sigma \theta(p) \, f(p, T_p\Sigma) \, d\mathcal{H}^n(p) \quad \text{for all } f \in C_c(\mathrm{Gr}_n(M)).
\]
We call $\theta$ the \emph{multiplicity function} of $V$.
\end{definition}

\begin{definition}[Integral varifold~{\cite[\S15, p.\,77]{Sim84}}]\label{def:integral-varifold}
A rectifiable $n$-varifold $V = \underline{v}(\Sigma, \theta)$ on $(M^{n+1}, g)$ is \emph{integral} if the multiplicity function $\theta \colon \Sigma \to (0, \infty)$ satisfies $\theta(p) \in \mathbb{Z}_{>0}$ for $\mathcal{H}^n$-a.e.\ $p \in \Sigma$.
\end{definition}

\begin{definition}[Stationary and stable varifolds~{\cite[Definition~16.3]{Sim84}}]\label{def:p2-stationary}
A varifold $V \in \mathcal{V}_n(M)$ is \emph{stationary} if its first variation vanishes:
\[
\delta V(X) := \int_{\mathrm{Gr}_n(M)} \operatorname{div}_S X \, dV = 0 \quad \text{for every } X \in \mathfrak{X}(M).
\]
A stationary varifold $V$ is \emph{stable} if, in addition, the second variation $\delta^2 V(\varphi, \varphi) \geq 0$ for all compactly supported normal deformations $\varphi$ of $\reg\, V$.
\end{definition}

\begin{definition}[Tangent cone~{\cite[Definition~42.3]{Sim84}}]\label{def:p2-tangent-cone}
Let $V$ be a stationary $n$-varifold in $M$ and $Z \in \spt\|V\|$. For $\rho > 0$, the \emph{blow-up} of $V$ at $Z$ with scale $\rho$ is the varifold $(\eta_{Z,\rho})_\# V$, where $\eta_{Z,\rho}(x) = (x - Z)/\rho$ in normal coordinates centred at $Z$. A \emph{tangent cone} of $V$ at $Z$ is any varifold limit $C = \lim_{j \to \infty} (\eta_{Z,\rho_j})_\# V$ for some sequence $\rho_j \to 0$. By the monotonicity formula (Proposition~\ref{prop:p2-monotonicity} below), at least one such limit exists; it is a stationary cone in $\mathbb{R}^{n+1}$.
\end{definition}

We next collect several quantitative tools from geometric measure theory that will be used in the integrality and regularity arguments of Sections~\ref{sec:integrality} and~\ref{sec:regularity}.

\begin{definition}[Varifold convergence~{\cite[\S38, p.\,228]{Sim84}}]\label{def:p2-varifold-convergence}
A sequence of $n$-varifolds $\{V_i\}$ \emph{converges} to $V \in \mathcal{V}_n(M)$ if $V_i \to V$ as Radon measures on $\mathrm{Gr}_n(M)$, i.e., $\int f \, dV_i \to \int f \, dV$ for every $f \in C_c(\mathrm{Gr}_n(M))$.
\end{definition}

\begin{proposition}[Monotonicity formula~{\cite[Theorem~17.6]{Sim84}}]\label{prop:p2-monotonicity}
Let $V$ be a stationary $n$-varifold in $(M^{n+1}, g)$ and $p \in \spt\|V\|$. There exists a constant $C = C(M, g) > 0$ such that the modified density ratio $e^{Cr}\, \omega_n^{-1} r^{-n} \|V\|(B_r(p))$ is monotone non-decreasing in $r$ (for $r$ small enough that the exponential map at $p$ is a diffeomorphism). In particular, at least one tangent cone exists at every point of $\spt\|V\|$.
\end{proposition}

\begin{definition}[Density~{\cite[Corollary~17.8]{Sim84}}]\label{def:p2-density}
For a stationary $n$-varifold $V$ and $p \in M$, the \emph{density} of $V$ at $p$ is
\[
\Theta(\|V\|, p) := \lim_{r \to 0} \omega_n^{-1} r^{-n} \|V\|(B_r(p)).
\]
By the monotonicity formula (Proposition~\ref{prop:p2-monotonicity}), this limit exists and $\Theta(\|V\|, p) \geq 0$ for every $p \in M$, with $\Theta(\|V\|, p) < \infty$ for every $p$ when $\|V\|(M) < \infty$.

If $V$ is integral, then $\Theta(\|V\|, p) \geq 1$ for every $p \in \spt\|V\|$.
\end{definition}

\begin{proposition}[Allard's rectifiability theorem~{\cite[Theorem~5.5(1)]{All72}; \cite[Theorem~42.4]{Sim84}}]\label{prop:allard-rectifiability}
Let $V$ be a stationary $n$-varifold in $(M^{n+1}, g)$ with $\Theta(\|V\|, p) > 0$ for $\|V\|$-a.e.\ $p$. Then $V$ is rectifiable: $V = \underline{v}(\Sigma, \theta)$ where $\Sigma \subset \spt\|V\|$ is $\mathcal{H}^n$-rectifiable and $\theta(p) = \Theta(\|V\|, p) > 0$ for $\mathcal{H}^n$-a.e.\ $p \in \Sigma$.
\end{proposition}

\begin{theorem}[Lebesgue--Besicovitch differentiation~{\cite[Theorem~5.8 and \S5.3]{Mag12}}]\label{thm:lebesgue-besicovitch}
Let $\mu$ be a locally finite Borel measure on $\mathbb{R}^{n+1}$ and $f \in L^1_{\mathrm{loc}}(\mu)$. Then $\mu$-a.e.\ $p$ is a \emph{Lebesgue point} of $f$:
\[
    \frac{1}{\mu(B_r(p))}\int_{B_r(p)} f\,d\mu \to f(p) \quad \text{as } r \to 0.
\]
The same holds on a Riemannian manifold $(M, g)$ for geodesic balls, since the statement is local and the metric is uniformly comparable to the Euclidean metric in normal coordinates.
\end{theorem}

\begin{definition}[Tilt-excess~{\cite[\S22, p.\,111]{Sim84}}]\label{def:p2-tilt-excess}
The \emph{tilt-excess} of a varifold $W \in \mathcal{V}_n(M)$ in $B_r(p)$ relative to an $n$-plane $P$ is
\[
E(W, B_r(p), P) := r^{-n} \int_{B_r(p) \times G(n,n+1)} |S - P|^2 \, dW(x, S).
\]
Intuitively, the tilt-excess is a scale-invariant $L^2$ measure of how much the tangent planes of $W$ deviate from the reference plane~$P$ inside $B_r(p)$. In the codimension-one case, when $V = \underline{v}(\Sigma, \theta)$ is rectifiable and $\Sigma$ is locally the graph of a function $u$ over $P$, the tilt-excess reduces (up to constants) to the normalised Dirichlet energy $r^{-n} \int |\nabla u|^2$; in particular, $E = 0$ if and only if $\Sigma \cap B_r(p)$ is a portion of a plane parallel to $P$.

\end{definition}

\begin{lemma}[Vanishing of tilt-excess]\label{lem:excess-vanishing}
Let $V = \underline{v}(\Sigma, \theta)$ be a rectifiable $n$-varifold on $(M^{n+1}, g)$. For $\mathcal{H}^n$-a.e.\ $p \in \Sigma$,
\[
    E(V, B_r(p), P) \to 0 \quad \text{as } r \to 0,
\]
where $P = T_p\Sigma$ is the approximate tangent plane.
\end{lemma}

\begin{proof}
Write $\mu = \|V\|$ and $f(x) = |T_x\Sigma - P|^2$, so that $E(V, B_r(p), P) = r^{-n}\int_{B_r(p)} f\,d\mu$. The function $f$ is bounded and $\mu$-measurable, with $f(p) = 0$. By the Lebesgue--Besicovitch differentiation theorem (Theorem~\ref{thm:lebesgue-besicovitch}), $\mu$-a.e.\ $p$ is a Lebesgue point of $f$:
\[
    \frac{1}{\mu(B_r(p))}\int_{B_r(p)} f\,d\mu \to f(p) = 0.
\]
Since $\mu(B_r(p)) = \|V\|(B_r(p)) \sim \theta(p)\,\omega_n\, r^n$ by the definition of density, we conclude $E(V, B_r(p), P) \to 0$.
\end{proof}

Finally, the following interpolation lemma from~\cite{CLS22} will be used to construct flat-continuous families connecting two nested Caccioppoli sets with controlled perimeter. Recall that a family $\{\partial\Omega_t\}_{t \in [0,1]}$ is \emph{$\mathcal{F}$-continuous} if $t \mapsto \Omega_t$ is continuous with respect to the flat norm $\mathcal{F}(\partial\Omega_s, \partial\Omega_t) = \mathrm{Vol}(\Omega_s \triangle \Omega_t)$.

\begin{lemma}[Interpolation lemma~{\cite[Lemma~1.12]{CLS22}}]\label{lem:interpolation}
Fix $L > 0$. For every $\varepsilon > 0$ there exists $\delta > 0$ such that the following holds. If $\Omega_0, \Omega_1$ are two Caccioppoli sets with $\Omega_0 \subset \Omega_1$, $\mathrm{Per}(\Omega_i) \leq L$ for $i = 0, 1$, and $\mathrm{Vol}(\Omega_1 \setminus \Omega_0) \leq \delta$, then there exists a nested $\mathcal{F}$-continuous family $\{\partial\Omega_t\}_{t \in [0,1]}$ with
\[
    \mathrm{Per}(\Omega_t) \leq \max\bigl\{\mathrm{Per}(\Omega_0),\, \mathrm{Per}(\Omega_1)\bigr\} + \varepsilon
\]
for all $t \in [0,1]$.
\end{lemma}

